% --- Archon physics formalization source begin ---
% archon:physics
% archon:covers PhyXMiniProblems/problem_phyx_mini_0998.lean
% archon:source-report reports/phyx_mini/problem_phyx_mini_0998.source.json

\chapter{Physics problem phyx\_mini\_0998}
\label{ch:PhyXMiniProblems_problem_phyx_mini_0998}

\paragraph{Problem source.}
\#\# Physical scenario

A disk of radius \textbackslash{}( 0.10 \textbackslash{}, \textbackslash{}text\{m\} \textbackslash{}) is oriented with its normal unit vector \textbackslash{}( \textbackslash{}hat\{\textbackslash{}mathbf\{n\}\} \textbackslash{}) at \textbackslash{}( 30\textasciicircum{}\textbackslash{}circ \textbackslash{}) to a uniform electric field \textbackslash{}( \textbackslash{}vec\{\textbackslash{}mathbf\{E\}\} \textbackslash{}) of magnitude \textbackslash{}( 2.0 \textbackslash{}times 10\textasciicircum{}3 \textbackslash{}, \textbackslash{}text\{N/C\} \textbackslash{}). (Since this isn't a closed surface, it has no “inside” or “outside.” That's why we have to specify the direction of \textbackslash{}( \textbackslash{}hat\{\textbackslash{}mathbf\{n\}\} \textbackslash{}) in the figure.)

\#\# Question

What is the electric flux through the disk?

\#\# Answer choices

- A: A: 108 \textbackslash{}

- B: B: 54 \textbackslash{}

- C: C: 58 \textbackslash{}

- D: D: 60 \textbackslash{}

\#\# Figure caption (auxiliary; use the image as primary evidence)

The image is a diagram featuring a circular disk or surface within a uniform electric field. Here's a detailed description of the image content:

- **Circular Surface**: A light-colored, circular disk is depicted at an angle, indicating a 3D view.
- **Radius Information**: The radius of the disk is labeled as \textbackslash{}( r = 0.10 \textbackslash{}, \textbackslash{}text\{m\} \textbackslash{}).
- **Normal Vector (\textbackslash{}(\textbackslash{}mathbf\{\textbackslash{}hat\{n\}\}\textbackslash{}))**: A black arrow represents the normal vector (\textbackslash{}(\textbackslash{}mathbf\{\textbackslash{}hat\{n\}\}\textbackslash{})) to the surface. The vector extends outward from the disk, forming an angle with the surface.
- **Angle**: The angle between the normal vector and the vertical is marked as \textbackslash{}(30\textasciicircum{}\textbackslash{}circ\textbackslash{}).
- **Electric Field (\textbackslash{}(\textbackslash{}mathbf\{E\}\textbackslash{}))**: Parallel pink arrows point horizontally across the image from left to right, representing a uniform electric field (\textbackslash{}(\textbackslash{}mathbf\{E\}\textbackslash{})) passing through the disk.
- **Right Triangle**: A small right triangle is depicted on the surface to emphasize the angle between the normal vector and the surface.

The illustration represents a typical physics scenario for analyzing electric flux through a surface at an angle within an electric field.

\#\# Dataset metadata

Electrostatics; Physical Model Grounding Reasoning, Spatial Relation Reasoning

\paragraph{Recorded answer/context.}
B: B: 54 \textbackslash{}

\paragraph{Figure/image path.}
/root/proposal\_for\_physic/science-mango/phyx\_mini\_run/phyx\_data/test\_image/998.png

\paragraph{Formalization target.}
create a compiling Lean file with sorry bodies at `PhyXMiniProblems/problem\_phyx\_mini\_0998.lean`. The Lean declarations must preserve the physical quantities, dimensions or dimensional roles, figure labels, governing-law hypotheses, and final relation expressed by this problem.
Use Mathlib/Physlib names found through LeanExplore where available. If a physics API is missing, introduce faithful local abstractions rather than scalar placeholder aliases.

\begin{theorem}[Physics formalization target]
\label{thm:physics:phyx_mini_0998:target}
\lean{PhyXMiniProblems.ProblemPhyXMini0998.problem_phyx_mini_0998}
\uses{def:physics:phyx-mini-0998:phyxminiproblems-problemphyxmini0998-electricfluxinnewtonsquaremeterspercoulomb, def:physics:phyx-mini-0998:phyxminiproblems-problemphyxmini0998-inclineddiskelectricfluxsetup, def:physics:phyx-mini-0998:phyxminiproblems-problemphyxmini0998-matchesproblemandprimaryfigure, def:physics:phyx-mini-0998:phyxminiproblems-problemphyxmini0998-hasphysicalcirculardiskgeometry, def:physics:phyx-mini-0998:phyxminiproblems-problemphyxmini0998-satisfiesfigurederivedorientation, def:physics:phyx-mini-0998:phyxminiproblems-problemphyxmini0998-satisfiesuniformelectricfieldmodel, def:physics:phyx-mini-0998:phyxminiproblems-problemphyxmini0998-satisfiesplanarelectricfluxlaw, def:physics:phyx-mini-0998:phyxminiproblems-problemphyxmini0998-recordeddatasetanswer, def:physics:phyx-mini-0998:phyxminiproblems-problemphyxmini0998-isuniqueclosestdisplayedchoice}
The assigned autoformalize agent should translate this physics problem into Lean declarations in the covered file, with theorem and lemma proof bodies written as `by sorry`.
\end{theorem}
\begin{proof}
This is an autoformalization task, not a proof task. Produce statements that can later be proved by the physics prover without weakening the physical model.
\end{proof}

% --- Archon declaration topology begin ---
\section{Declaration topology}
\begin{definition}[electric Field Strength Dimension]
  \label{def:physics:phyx-mini-0998:phyxminiproblems-problemphyxmini0998-electricfieldstrengthdimension}
  \lean{PhyXMiniProblems.ProblemPhyXMini0998.electricFieldStrengthDimension}
  Electric-field strength has dimension M L T⁻² C⁻¹ (N/C)
\end{definition}
\begin{proof}
  This is the declared model definition of electric Field Strength Dimension.
\end{proof}
\begin{definition}[electric Flux Dimension]
  \label{def:physics:phyx-mini-0998:phyxminiproblems-problemphyxmini0998-electricfluxdimension}
  \lean{PhyXMiniProblems.ProblemPhyXMini0998.electricFluxDimension}
  \uses{def:physics:phyx-mini-0998:phyxminiproblems-problemphyxmini0998-electricfieldstrengthdimension}
  Electric flux has dimension M L³ T⁻² C⁻¹ (N m²/C)
\end{definition}
\begin{proof}
  This is the declared model definition of electric Flux Dimension.
\end{proof}
\begin{definition}[Length Quantity]
  \label{def:physics:phyx-mini-0998:phyxminiproblems-problemphyxmini0998-lengthquantity}
  \lean{PhyXMiniProblems.ProblemPhyXMini0998.LengthQuantity}
  A nonnegative, unit-independent physical length
\end{definition}
\begin{proof}
  This is the declared model definition of Length Quantity.
\end{proof}
\begin{definition}[Area Quantity]
  \label{def:physics:phyx-mini-0998:phyxminiproblems-problemphyxmini0998-areaquantity}
  \lean{PhyXMiniProblems.ProblemPhyXMini0998.AreaQuantity}
  A nonnegative, unit-independent physical area
\end{definition}
\begin{proof}
  This is the declared model definition of Area Quantity.
\end{proof}
\begin{definition}[Electric Field Strength Magnitude]
  \label{def:physics:phyx-mini-0998:phyxminiproblems-problemphyxmini0998-electricfieldstrengthmagnitude}
  \lean{PhyXMiniProblems.ProblemPhyXMini0998.ElectricFieldStrengthMagnitude}
  \uses{def:physics:phyx-mini-0998:phyxminiproblems-problemphyxmini0998-electricfieldstrengthdimension}
  A nonnegative, unit-independent electric-field-strength magnitude
\end{definition}
\begin{proof}
  This is the declared model definition of Electric Field Strength Magnitude.
\end{proof}
\begin{definition}[Electric Flux Quantity]
  \label{def:physics:phyx-mini-0998:phyxminiproblems-problemphyxmini0998-electricfluxquantity}
  \lean{PhyXMiniProblems.ProblemPhyXMini0998.ElectricFluxQuantity}
  \uses{def:physics:phyx-mini-0998:phyxminiproblems-problemphyxmini0998-electricfluxdimension}
  Signed electric flux through an oriented surface
\end{definition}
\begin{proof}
  This is the declared model definition of Electric Flux Quantity.
\end{proof}
\begin{definition}[length Readout]
  \label{def:physics:phyx-mini-0998:phyxminiproblems-problemphyxmini0998-lengthreadout}
  \lean{PhyXMiniProblems.ProblemPhyXMini0998.lengthReadout}
  \uses{def:physics:phyx-mini-0998:phyxminiproblems-problemphyxmini0998-lengthquantity}
  Read a physical length in a selected named length unit
\end{definition}
\begin{proof}
  This is the declared model definition of length Readout.
\end{proof}
\begin{definition}[area Readout]
  \label{def:physics:phyx-mini-0998:phyxminiproblems-problemphyxmini0998-areareadout}
  \lean{PhyXMiniProblems.ProblemPhyXMini0998.areaReadout}
  \uses{def:physics:phyx-mini-0998:phyxminiproblems-problemphyxmini0998-areaquantity}
  Read a physical area in the square of a selected named length unit
\end{definition}
\begin{proof}
  This is the declared model definition of area Readout.
\end{proof}
\begin{definition}[length In Meters]
  \label{def:physics:phyx-mini-0998:phyxminiproblems-problemphyxmini0998-lengthinmeters}
  \lean{PhyXMiniProblems.ProblemPhyXMini0998.lengthInMeters}
  \uses{def:physics:phyx-mini-0998:phyxminiproblems-problemphyxmini0998-lengthquantity, def:physics:phyx-mini-0998:phyxminiproblems-problemphyxmini0998-lengthreadout}
  Read a physical length in metres
\end{definition}
\begin{proof}
  This is the declared model definition of length In Meters.
\end{proof}
\begin{definition}[area In Square Meters]
  \label{def:physics:phyx-mini-0998:phyxminiproblems-problemphyxmini0998-areainsquaremeters}
  \lean{PhyXMiniProblems.ProblemPhyXMini0998.areaInSquareMeters}
  \uses{def:physics:phyx-mini-0998:phyxminiproblems-problemphyxmini0998-areaquantity, def:physics:phyx-mini-0998:phyxminiproblems-problemphyxmini0998-areareadout}
  Read a physical area in square metres
\end{definition}
\begin{proof}
  This is the declared model definition of area In Square Meters.
\end{proof}
\begin{definition}[electric Field Strength In Newtons Per Coulomb]
  \label{def:physics:phyx-mini-0998:phyxminiproblems-problemphyxmini0998-electricfieldstrengthinnewtonspercoulomb}
  \lean{PhyXMiniProblems.ProblemPhyXMini0998.electricFieldStrengthInNewtonsPerCoulomb}
  \uses{def:physics:phyx-mini-0998:phyxminiproblems-problemphyxmini0998-electricfieldstrengthmagnitude}
  Read an electric-field-strength magnitude in newtons per coulomb
\end{definition}
\begin{proof}
  This is the declared model definition of electric Field Strength In Newtons Per Coulomb.
\end{proof}
\begin{definition}[electric Flux In Newton Square Meters Per Coulomb]
  \label{def:physics:phyx-mini-0998:phyxminiproblems-problemphyxmini0998-electricfluxinnewtonsquaremeterspercoulomb}
  \lean{PhyXMiniProblems.ProblemPhyXMini0998.electricFluxInNewtonSquareMetersPerCoulomb}
  \uses{def:physics:phyx-mini-0998:phyxminiproblems-problemphyxmini0998-electricfluxquantity}
  Read signed electric flux in newton-square-metres per coulomb
\end{definition}
\begin{proof}
  This is the declared model definition of electric Flux In Newton Square Meters Per Coulomb.
\end{proof}
\begin{definition}[degrees To Radians]
  \label{def:physics:phyx-mini-0998:phyxminiproblems-problemphyxmini0998-degreestoradians}
  \lean{PhyXMiniProblems.ProblemPhyXMini0998.degreesToRadians}
  Convert a degree readout to a real angle in radians
\end{definition}
\begin{proof}
  This is the declared model definition of degrees To Radians.
\end{proof}
\begin{definition}[Diagram Arrow Direction]
  \label{def:physics:phyx-mini-0998:phyxminiproblems-problemphyxmini0998-diagramarrowdirection}
  \lean{PhyXMiniProblems.ProblemPhyXMini0998.DiagramArrowDirection}
  Qualitative arrow directions visible in the supplied two-dimensional view
\end{definition}
\begin{proof}
  This is the declared model definition of Diagram Arrow Direction.
\end{proof}
\begin{definition}[Figure Feature]
  \label{def:physics:phyx-mini-0998:phyxminiproblems-problemphyxmini0998-figurefeature}
  \lean{PhyXMiniProblems.ProblemPhyXMini0998.FigureFeature}
  Named visual features visible in phyx\_data/test\_image/998.png
\end{definition}
\begin{proof}
  This is the declared model definition of Figure Feature.
\end{proof}
\begin{definition}[Inclined Disk Figure]
  \label{def:physics:phyx-mini-0998:phyxminiproblems-problemphyxmini0998-inclineddiskfigure}
  \lean{PhyXMiniProblems.ProblemPhyXMini0998.InclinedDiskFigure}
  \uses{def:physics:phyx-mini-0998:phyxminiproblems-problemphyxmini0998-diagramarrowdirection, def:physics:phyx-mini-0998:phyxminiproblems-problemphyxmini0998-figurefeature}
  Literal problem and raster data. No flux value or answer label is stored in this record
\end{definition}
\begin{proof}
  This is the declared model definition of Inclined Disk Figure.
\end{proof}
\begin{definition}[Oriented Circular Disk]
  \label{def:physics:phyx-mini-0998:phyxminiproblems-problemphyxmini0998-orientedcirculardisk}
  \lean{PhyXMiniProblems.ProblemPhyXMini0998.OrientedCircularDisk}
  \uses{def:physics:phyx-mini-0998:phyxminiproblems-problemphyxmini0998-lengthquantity, def:physics:phyx-mini-0998:phyxminiproblems-problemphyxmini0998-areaquantity}
  An oriented circular disk in three-dimensional physical space. The set and representative point preserve the surface's geometric role; its radius and area remain dimensionful quantities
\end{definition}
\begin{proof}
  This is the declared model definition of Oriented Circular Disk.
\end{proof}
\begin{definition}[Inclined Disk Electric Flux Setup]
  \label{def:physics:phyx-mini-0998:phyxminiproblems-problemphyxmini0998-inclineddiskelectricfluxsetup}
  \lean{PhyXMiniProblems.ProblemPhyXMini0998.InclinedDiskElectricFluxSetup}
  \uses{def:physics:phyx-mini-0998:phyxminiproblems-problemphyxmini0998-electricfieldstrengthmagnitude, def:physics:phyx-mini-0998:phyxminiproblems-problemphyxmini0998-electricfluxquantity, def:physics:phyx-mini-0998:phyxminiproblems-problemphyxmini0998-inclineddiskfigure, def:physics:phyx-mini-0998:phyxminiproblems-problemphyxmini0998-orientedcirculardisk}
  The independent physical quantities and observables. In particular, electricFlux is not defined from an answer choice or from the requested numerical value
\end{definition}
\begin{proof}
  This is the declared model definition of Inclined Disk Electric Flux Setup.
\end{proof}
\begin{definition}[Matches Problem And Primary Figure]
  \label{def:physics:phyx-mini-0998:phyxminiproblems-problemphyxmini0998-matchesproblemandprimaryfigure}
  \lean{PhyXMiniProblems.ProblemPhyXMini0998.MatchesProblemAndPrimaryFigure}
  \uses{def:physics:phyx-mini-0998:phyxminiproblems-problemphyxmini0998-lengthinmeters, def:physics:phyx-mini-0998:phyxminiproblems-problemphyxmini0998-electricfieldstrengthinnewtonspercoulomb, def:physics:phyx-mini-0998:phyxminiproblems-problemphyxmini0998-inclineddiskelectricfluxsetup}
  Direct source and primary-image readouts. The prose supplies the field magnitude, while the raster supplies the radius label, arrow directions, chosen normal, and acute 30° marker. No flux result occurs here
\end{definition}
\begin{proof}
  This is the declared model definition of Matches Problem And Primary Figure.
\end{proof}
\begin{definition}[Has Physical Circular Disk Geometry]
  \label{def:physics:phyx-mini-0998:phyxminiproblems-problemphyxmini0998-hasphysicalcirculardiskgeometry}
  \lean{PhyXMiniProblems.ProblemPhyXMini0998.HasPhysicalCircularDiskGeometry}
  \uses{def:physics:phyx-mini-0998:phyxminiproblems-problemphyxmini0998-lengthinmeters, def:physics:phyx-mini-0998:phyxminiproblems-problemphyxmini0998-areainsquaremeters, def:physics:phyx-mini-0998:phyxminiproblems-problemphyxmini0998-inclineddiskelectricfluxsetup}
  Regularity and circular-area geometry of the open disk. The area relation is a general geometric law and contains no flux value
\end{definition}
\begin{proof}
  This is the declared model definition of Has Physical Circular Disk Geometry.
\end{proof}
\begin{definition}[Satisfies Figure Derived Orientation]
  \label{def:physics:phyx-mini-0998:phyxminiproblems-problemphyxmini0998-satisfiesfigurederivedorientation}
  \lean{PhyXMiniProblems.ProblemPhyXMini0998.SatisfiesFigureDerivedOrientation}
  \uses{def:physics:phyx-mini-0998:phyxminiproblems-problemphyxmini0998-degreestoradians, def:physics:phyx-mini-0998:phyxminiproblems-problemphyxmini0998-inclineddiskelectricfluxsetup}
  The figure-derived spatial relation: the field direction meets the chosen oriented unit normal at the marked acute angle
\end{definition}
\begin{proof}
  This is the declared model definition of Satisfies Figure Derived Orientation.
\end{proof}
\begin{definition}[Satisfies Uniform Electric Field Model]
  \label{def:physics:phyx-mini-0998:phyxminiproblems-problemphyxmini0998-satisfiesuniformelectricfieldmodel}
  \lean{PhyXMiniProblems.ProblemPhyXMini0998.SatisfiesUniformElectricFieldModel}
  \uses{def:physics:phyx-mini-0998:phyxminiproblems-problemphyxmini0998-electricfieldstrengthinnewtonspercoulomb, def:physics:phyx-mini-0998:phyxminiproblems-problemphyxmini0998-inclineddiskelectricfluxsetup}
  Uniform-field idealization on the observation time slice. Physlib's vector field is calibrated in coherent-SI components by the independent dimensionful strength magnitude and unit direction. This premise mentions no flux
\end{definition}
\begin{proof}
  This is the declared model definition of Satisfies Uniform Electric Field Model.
\end{proof}
\begin{definition}[Satisfies Planar Electric Flux Law]
  \label{def:physics:phyx-mini-0998:phyxminiproblems-problemphyxmini0998-satisfiesplanarelectricfluxlaw}
  \lean{PhyXMiniProblems.ProblemPhyXMini0998.SatisfiesPlanarElectricFluxLaw}
  \uses{def:physics:phyx-mini-0998:phyxminiproblems-problemphyxmini0998-areainsquaremeters, def:physics:phyx-mini-0998:phyxminiproblems-problemphyxmini0998-electricfluxinnewtonsquaremeterspercoulomb, def:physics:phyx-mini-0998:phyxminiproblems-problemphyxmini0998-inclineddiskelectricfluxsetup}
  For an oriented planar surface in a uniform field, electric flux is A * (E dot n). This is the governing surface-flux law, not the requested numerical evaluation
\end{definition}
\begin{proof}
  This is the declared model definition of Satisfies Planar Electric Flux Law.
\end{proof}
\begin{lemma}[inclined Disk Electric Flux exact]
  \label{lem:physics:phyx-mini-0998:phyxminiproblems-problemphyxmini0998-inclineddiskelectricflux-exact}
  \lean{PhyXMiniProblems.ProblemPhyXMini0998.inclinedDiskElectricFlux_exact}
  \uses{def:physics:phyx-mini-0998:phyxminiproblems-problemphyxmini0998-electricfluxinnewtonsquaremeterspercoulomb, def:physics:phyx-mini-0998:phyxminiproblems-problemphyxmini0998-inclineddiskelectricfluxsetup, def:physics:phyx-mini-0998:phyxminiproblems-problemphyxmini0998-matchesproblemandprimaryfigure, def:physics:phyx-mini-0998:phyxminiproblems-problemphyxmini0998-hasphysicalcirculardiskgeometry, def:physics:phyx-mini-0998:phyxminiproblems-problemphyxmini0998-satisfiesfigurederivedorientation, def:physics:phyx-mini-0998:phyxminiproblems-problemphyxmini0998-satisfiesuniformelectricfieldmodel, def:physics:phyx-mini-0998:phyxminiproblems-problemphyxmini0998-satisfiesplanarelectricfluxlaw}
  Substituting r = 0.10 m, E = 2000 N/C, and theta = 30° into Phi = E * pi * r\textasciicircum{}2 * cos theta gives the exact idealized flux 10 * pi * sqrt 3 N m\textasciicircum{}2/C
\end{lemma}
\begin{proof}
  The conclusion follows from the governing relations and figure readouts named in the statement of inclined Disk Electric Flux exact.
\end{proof}
\begin{definition}[Answer Choice]
  \label{def:physics:phyx-mini-0998:phyxminiproblems-problemphyxmini0998-answerchoice}
  \lean{PhyXMiniProblems.ProblemPhyXMini0998.AnswerChoice}
  Labels of the four electric-flux choices printed with the question
\end{definition}
\begin{proof}
  This is the declared model definition of Answer Choice.
\end{proof}
\begin{definition}[displayed Electric Flux]
  \label{def:physics:phyx-mini-0998:phyxminiproblems-problemphyxmini0998-displayedelectricflux}
  \lean{PhyXMiniProblems.ProblemPhyXMini0998.displayedElectricFlux}
  \uses{def:physics:phyx-mini-0998:phyxminiproblems-problemphyxmini0998-answerchoice}
  Electric-flux readout in N m²/C printed beside an answer choice
\end{definition}
\begin{proof}
  This is the declared model definition of displayed Electric Flux.
\end{proof}
\begin{definition}[recorded Dataset Answer]
  \label{def:physics:phyx-mini-0998:phyxminiproblems-problemphyxmini0998-recordeddatasetanswer}
  \lean{PhyXMiniProblems.ProblemPhyXMini0998.recordedDatasetAnswer}
  \uses{def:physics:phyx-mini-0998:phyxminiproblems-problemphyxmini0998-answerchoice}
  The answer label recorded by the source dataset, retained as metadata
\end{definition}
\begin{proof}
  This is the declared model definition of recorded Dataset Answer.
\end{proof}
\begin{definition}[Is Unique Closest Displayed Choice]
  \label{def:physics:phyx-mini-0998:phyxminiproblems-problemphyxmini0998-isuniqueclosestdisplayedchoice}
  \lean{PhyXMiniProblems.ProblemPhyXMini0998.IsUniqueClosestDisplayedChoice}
  \uses{def:physics:phyx-mini-0998:phyxminiproblems-problemphyxmini0998-electricfluxinnewtonsquaremeterspercoulomb, def:physics:phyx-mini-0998:phyxminiproblems-problemphyxmini0998-inclineddiskelectricfluxsetup, def:physics:phyx-mini-0998:phyxminiproblems-problemphyxmini0998-answerchoice, def:physics:phyx-mini-0998:phyxminiproblems-problemphyxmini0998-displayedelectricflux}
  A choice is uniquely closest to the idealized signed flux when its absolute error is strictly less than that of every other displayed readout. This definition does not privilege a particular label
\end{definition}
\begin{proof}
  This is the declared model definition of Is Unique Closest Displayed Choice.
\end{proof}
% --- Archon declaration topology end ---
% --- Archon physics formalization source end ---
